\section{Zeros, conjugated dynamics, and the scope of recovery}
\label{sec:consequences}

\subsection{Exact zero preservation}

For fixed $\tau\ge0$, define
\begin{equation}
 m_\tau(z)=q(\tau)\ee^{-\tau z^2}.
 \label{eq:multiplier}
\end{equation}
It is entire and nowhere zero. Hence if, for some fixed $\lambda$,
$H_\lambda(z)=(z-z_0)^r h(z)$ with $h(z_0)\ne0$, then
\begin{equation*}
 \Obs_{\lambda,\tau}(z)=(z-z_0)^r m_\tau(z)h(z)
\end{equation*}
has exactly the same multiplicity $r$ at $z_0$. Since $H_\lambda(0)>0$,
the identically zero case does not occur. Thus zeros, multiplicities,
and their collision loci as $\lambda$ varies are preserved exactly.
This conclusion follows from invertible multiplication; it supplies
no new inequality constraining where those collisions can occur.

\subsection{The conjugated Newman equation}

Introduce the differential operator
\begin{equation}
 \mathscr D_\tau=\partial_z+2\tau z.
 \label{eq:covariant-derivative}
\end{equation}
For any differentiable $h$,
\begin{equation}
 \mathscr D_\tau(m_\tau h)=m_\tau h'.
 \label{eq:conjugation}
\end{equation}
Iterating this identity and using \eqref{eq:newman-heat} gives
\begin{equation}
 \boxed{\partial_\lambda\Obs_{\lambda,\tau}
       =-\mathscr D_\tau^2\Obs_{\lambda,\tau}}
 \label{eq:conjugated-heat}
\end{equation}
or, in ordinary derivatives,
\begin{equation}
 \partial_\lambda\Obs
 =-\partial_z^2\Obs-4\tau z\partial_z\Obs
                 -(2\tau+4\tau^2z^2)\Obs.
 \label{eq:conjugated-heat-expanded}
\end{equation}
The term $2\tau\Obs$ comes from differentiating the coefficient
$2\tau z$ inside $\mathscr D_\tau^2$ and must be retained.
For every finite $a$,
$\partial_\lambda^a\Obs=(-1)^a\mathscr D_\tau^{2a}\Obs$;
these are finite differential identities justified by the preceding
regularity, not claims about infinite series of operators.

The proper-time equation is instead
\begin{equation}
 \partial_\tau\Obs_{\lambda,\tau}(z)
 =\left(\frac{q'(\tau)}{q(\tau)}-z^2\right)
                         \Obs_{\lambda,\tau}(z).
 \label{eq:proper-time-response}
\end{equation}
Equations~\eqref{eq:conjugated-heat} and
\eqref{eq:proper-time-response} show directly why Newman deformation
and auxiliary proper-time evolution are different operations.

\subsection{Inversion and the first derivative pair}

The exact inverse is
\begin{equation}
 H_\lambda(z)=q(\tau)^{-1}\ee^{\tau z^2}
                         \Obs_{\lambda,\tau}(z).
 \label{eq:inverse}
\end{equation}
For real $k$ and $\tau>0$, its absolute multiplier is
$q(\tau)^{-1}\ee^{\tau k^2}$, which tends to infinity as
$|k|\to\infty$. Hence the construction does not give an inverse
bounded uniformly over all real frequencies. On a fixed window
$|k|\le K$, $0\le\tau\le T$, the scalar inverse is bounded by
$q(T)^{-1}\ee^{TK^2}$ because $q$ is positive and decreasing.
This is a finite-window attenuation calculation, not a source-specific
collision estimate.

The pair of an amplitude and its first derivative is the natural next
diagnostic. Fix $s>0$ to specify the relative units of derivative error
and value error, and define, at a real frequency,
\begin{equation}
 X=\begin{pmatrix}H_\lambda(k)\\sH_\lambda'(k)\end{pmatrix},
 \qquad
 Y=\begin{pmatrix}\Obs_{\lambda,\tau}(k)\\
                         s\partial_k\Obs_{\lambda,\tau}(k)\end{pmatrix}.
\end{equation}
Differentiating \eqref{eq:response} once gives the exact algebraic map
\begin{equation}
 Y=mP_aX,\qquad m=q(\tau)\ee^{-\tau k^2},\qquad
 a=2s\tau k,\qquad
 P_a=\begin{pmatrix}1&0\\-a&1\end{pmatrix}.
 \label{eq:first-jet}
\end{equation}
Equivalently,
\begin{equation}
 P_a^{-1}Y=
 \begin{pmatrix}\Obs\\s(\partial_k+2\tau k)\Obs\end{pmatrix}
 =mX.
 \label{eq:covariant-pair}
\end{equation}
This diagonal form does not remove measurement error by itself. Under
the same change of coordinates an additive error set $E$ becomes
$P_a^{-1}E$. Similarly, source uncertainty $E$ propagated through the
exact map $\mathsf A=mP_a$ and then inverted returns
$\mathsf A^{-1}(\mathsf A E)=E$.
Changing the error model after changing coordinates would compare
different information. These elementary observations motivate the
pending sharp conditioning task in Section~\ref{sec:outstanding}; no
sharp first-jet recovery theorem is claimed here.

\subsection{What the representation establishes}

The construction proves a complete analytic realization of an exact
response under specified preparation and observation laws. It retains
the geometric spectrum, every physical angular shell, the full radial
integral, the passive reference, and the absolute gain. Its validity
does not rest on numerical truncation. It does not derive the source
scan or imprint from free cylinder dynamics, nor does it derive the
linear phase-sensitive detector from an intrinsic measurement model.

The positivity of $L$, $R$, and $\Phi$ also has distinct meanings.
The Hamiltonian is bounded below, the bank is a positive normalized
angular operator, and the source is positive. A real-frequency
observation nevertheless integrates the source against an oscillatory
cosine. Its zeros describe cancellation in this amplitude, not a zero
eigenvalue of the positive Hamiltonian. Hilbert-space positivity alone
therefore does not exclude a multiple zero of $H_\lambda$.

The existing control families make this distinction concrete
\cite{CollisionControls}. Their role is to test the premises of a
proposed estimate, rather than to replace the canonical source:
\begin{itemize}
\item For applicable Euler-factor controls, smooth positive-source
Fourier and heat operations persist, while the unchanged canonical
theta arithmetic fails. A generic representation of this kind rejects
none of their recorded collisions.
\item Weighted-square and finite-prime constructions can retain
positive half-line formulas while losing complete unit coefficients
and the canonical reflection identity. Matching finitely many endpoint
conditions does not establish the global smooth domains used here.
\item Applicable reflected super-Gaussian or smooth positive primitive
controls retain smoothness, positivity, reflection, and exponential
moments without the exact theta expansion. Where the analytic
hypotheses hold, the same nonvanishing multiplier preserves their
collisions as well.
\item Atomic or finite-profile controls can have Fourier observations
and moments without satisfying the smooth field-domain hypotheses.
No smooth-field realization for them follows from the present proof.
\end{itemize}
Consequently, naming the canonical source identifies the input but
does not yet quantitatively separate it from colliding controls.
Any future collision estimate must identify the additional arithmetic
or analytic property that makes such a separation possible.
